\section{Discussion}
\paragraph{What the agents did and did not do.} The models designed reasonable first crops and moved their final policy towards better ones over the year (Figure~\ref{fig:e1}, middle), but they did not use the facility as an instrument while it ran: no crop was ever stopped, the Full and Endpoint branches usually planted the same second wave, and process readings and the process predictor were consulted a handful of times per campaign. The E2 null therefore reflects how the agents used the interface rather than a lack of information in it; the baseline, which by construction reads the process predictor at mid-season, did not gain from it either, which suggests that the in-season signal available through the predictor added little in this task.

\paragraph{Gathering versus reading evidence.} E3 separates the two. A one-shot reading of the same final evidence produced worse recommendations than the campaign conversation that gathered it, while a fixed numerical Reader did at least as well as the gatherer on three of the four sources. Within the models tested, the limiting step was turning a year of noisy, confounded crop outcomes into a recommendation, not only collecting them.

\paragraph{What the baseline result means.} A simple trust-region search that starts from a standard grower setting and recommends only what it tried beat the models in the specified facility, yet captured only a small part of the headroom found by an offline numerical search. Its policies also differ in operating style (more light, less heat), and the comparison reverses when neighbouring spaces exchange more heat. SlowLab is therefore informative about how agents run asynchronous, irreversible experiments, and its E1 numbers should be read together with the sensitivity analysis rather than as a general ranking of methods for greenhouse control.

\section{Conclusion}
SlowLab tests whether an agent can run a year of slow, irreversible experiments in a shared facility and turn them into a good operating decision. Under a preregistered protocol, three low-cost language models learned within the year but stopped at the level of a standard grower setting, did not exploit in-season process measurements, and read final evidence worse than they gathered it. We will release the benchmark, its locked protocol and the public records of the run so that stronger agents can be measured against the same sites.
